\chapter{Architecture}

The firmware follows a deterministic execution cycle: initialization in \texttt{setup()} and periodic module updates in \texttt{loop()}. Each module exposes \texttt{begin()} and \texttt{update()} to separate configuration from runtime logic, reducing the risk of inconsistent states.

\section{Logical layout}
\begin{itemize}
  \item \textbf{Turnstile}: reads IN/OUT buttons, updates the people count, and moves the servo.
  \item \textbf{Stoplight}: displays availability using green/red LEDs.
  \item \textbf{Thermostat}: reads the NTC sensor and decides heating/cooling.
  \item \textbf{LightingControl}: reads the photoresistor and adjusts the dimmer.
  \item \textbf{HumidifierControl}: reads the DHT22 and turns humidity control on or off.
  \item \textbf{DisplayPanel}: alternates LCD pages with status, date/time, and diagnostics.
\end{itemize}

\section{Data flow}
The flow follows three constant steps:
\begin{enumerate}
  \item \textbf{Acquisition}: analog or digital readings from sensors.
  \item \textbf{Validation}: error checks and threshold filters.
  \item \textbf{Actuation}: LED, PWM, or servo based on decisions.
\end{enumerate}
The display aggregates module states and presents a concise view for the operator.

\section{Timing cycles}
Some modules use internal timing to avoid rapid transitions or flicker. In particular:
\begin{itemize}
  \item the thermostat applies a pause before switching mode;
  \item the display caps refresh at 1 second and switches page every 5 seconds.
\end{itemize}
These choices prevent oscillations and improve perceived stability.

\section{Key parameters}
Target values are configured in the main file:\newline
\texttt{20.0 C} for temperature, \texttt{65\%} for humidity, \texttt{200 lux} for lighting, \texttt{5 people} as the maximum capacity.
